% ===================================================================
%  ТАБЛИЦЯ № 8 — Жнива. — Полювання, Збір винограду, Риболовля.
%  Traduction 2026 d'apres texto/io, controlee sur texto/fr, texto/en
%  et texto/es. Consigne : voir texto/uk/10-tablycia-01.tex.
%
%  L'appel de note est dans le TITRE de la scene — « Жнива (*) » — et
%  non dans un alinea : c'est ainsi que l'ido le compose.
%
%  LE BLOC c2-03-1 EST UN ECART DECLARE dans renvoji.py, comme dans
%  toutes les colonnes : l'ordre des renvois y suit le francais et non
%  l'ido, et c'est voulu.
% ===================================================================

%%K t08-noto-1 noto
\VUnotes{143.2mm}{%
\textit{(*) Бо це слово неточне: жнива льону, збір винограду, збір яблук?
Можливо, краще було б узяти «} \VUgras{mesono} \textit{», запропоноване в}
Mondo \textit{XIV, стор. 242. Але поки цей новий корінь не прийнято Академією
і він чекає на обговорення за звичайними правилами ідо-руху.}%
}

%%K t08-apar-1 apar
\VUsaut{5.0mm}
\VUtitre{15.0pt}{60}{ТАБЛИЦЯ № 8}
\VUsaut{3.0mm}
\VUcentre{13.2pt}{90}{\textit{Перша сцена.}}
\VUsaut{3.0mm}
\VUpk{10.2pt}{40}{Жнива (*).}
\VUsaut{4.0mm}
\VUpk{10.2pt}{40}{Обличчя полів.}
\VUsaut{3.0mm}

%%K t08-c1-01-1 p
1. --- З приходом \VUgras{літа} поля знову перемінили обличчя. Насіння, ввірене
борозні, зійшло; із землі піднялася маленька зелена рослина і дала свій колос.
Зерно налилося, потім достигло, і тепер лишається тільки зібрати жниво.

%%K t08-c1-02-1 p
2. --- \VUgras{Польові квіти} розквітли. \VUgras{Волошки} \textsuperscript{(1)},
\VUgras{маки} \textsuperscript{(2)} і \VUgras{наперстянки}
\textsuperscript{(3)} вносять у рівний тон хлібів веселу суперечку своїх свіжих
\VUgras{віночків}.

%%K t08-c1-03-1 p
3. --- Плодові дерева вбралися в листя; плоди достигли. Маленька
\VUgras{вишня} \textsuperscript{(4)} гнеться під прекрасними \VUgras{вишнями}
\textsuperscript{(5)}, які діти рвуть і їдять із насолодою.

%%K t08-c1-04-1 p
4. --- Отара може тепер цілий день бути на волі.

%%K t08-c1-04-2 p
\VUgras{Ягнята} \textsuperscript{(6)} виросли і стали славними \VUgras{вівцями}
\textsuperscript{(7)} і славними \VUgras{баранами} \textsuperscript{(8)} з густим
руном. Вони мирно скубуть \VUgras{траву} на \VUgras{пасовиську}
\textsuperscript{(9)}, засіяному \VUgras{стокротками}.

%%K t08-c1-04-3 p
Скоро з цих тварин острижуть \VUgras{вовну}, з якої роблять сукно на наш одяг.

%%K t08-tit-2 sub
\VUsaut{5.0mm}
\VUpk{10.2pt}{40}{Пастух.}
\VUsaut{3.4mm}

%%K t08-c1-05-1 p
5. --- \VUgras{Пастух} \textsuperscript{(10)}, здається, не дуже-то на них
дивиться. Його займає тільки гроза, що проходить на обрії, а \VUgras{вівчар}
\textsuperscript{(13)}, підносячи до рота свою \VUgras{баклагу}
\textsuperscript{(14)}, здається ще безжурнішим.

%%K t08-c1-06-1 p
6. --- Спека тисне; ось і \VUgras{собака} \textsuperscript{(15)} іде освіжитися;
він хлебче свіжу воду \VUgras{струмка} \textsuperscript{(16)} і лякає бідну
маленьку \VUgras{жабу} \textsuperscript{(17)}, що причаїлася в траві на березі. Та
стрибає в прозору воду, де цвіте \VUgras{латаття} \textsuperscript{(18)}.

%%K t08-c1-07-1 p
7. --- \VUgras{Плиска} \textsuperscript{(19)} купається біля маленького
\VUgras{джерела} \textsuperscript{(20)}, що б'є між двома каменями і ховається під
\VUgras{колючим чагарником} \textsuperscript{(21)} і \VUgras{кущами глоду}
\textsuperscript{(22)}.

%%K t08-tit-3 sub
\VUsaut{5.0mm}
\VUpk{10.2pt}{40}{Жнива.}
\VUsaut{3.4mm}

%%K t08-c1-08-1 p
8. --- Для сільських дітей збирання хліба --- час дуже веселий. Вони весело
\VUgras{перекидаються} \textsuperscript{(24)} на \VUgras{копицях соломи}
\textsuperscript{(23)}, помагають \VUgras{женцям} \textsuperscript{(26)}, і поки
ті косять \VUgras{ниву} \textsuperscript{(27)} своїми \VUgras{косами}
\textsuperscript{(25)}, добре наклепаними на \VUgras{бруску}
\textsuperscript{(30)}, вони збирають хліб і в'яжуть його в \VUgras{снопи}
\textsuperscript{(31)} або \VUgras{в'язки} \textsuperscript{(32)}.

%%K t08-c1-09-1 p
9. --- Інколи найбільшим дають зрізати \VUgras{серпом} \textsuperscript{(12)}
\VUgras{золоті стебла} \textsuperscript{(28)}, що несуть важкі \VUgras{колоски}
\textsuperscript{(29)}, повні зерна; а маленькі, доглядаючи \VUgras{худобу}
\textsuperscript{(33)}, що пасеться на \VUgras{луці} \textsuperscript{(34)} в тіні
великої \VUgras{липи} \textsuperscript{(35)}, ловлять своїм \VUgras{сачком}
\textsuperscript{(36)} гарних \VUgras{метеликів}.

%%K t08-c1-09-2 p
Інші навіть вилазять на \VUgras{воза} \textsuperscript{(37)}, щоб укласти снопи,
які їм подають на \VUgras{вилах} \textsuperscript{(38)}.

%%K t08-c1-10-1 p
10. --- По всій \VUgras{рівнині} \textsuperscript{(39)}, де злітають
\VUgras{жайворонки} \textsuperscript{(41)}, іде велика метушня. Усі зайняті
жнивами. Але поки бідняки, як і давніше, беруться за старі знаряддя ---
\VUgras{коси, серпи} і \VUgras{ціпи} \textsuperscript{(40)} --- багаті
землевласники жнуть свою пшеницю або своє \VUgras{жито} \textsuperscript{(43)}
\VUgras{жниваркою} \textsuperscript{(42)}. Тоді, не возячи снопів до клуні, тут
же на місці складають великі \VUgras{скирти} \textsuperscript{(44)}.

%%K t08-c1-11-1 p
11. --- Потім підвозять \VUgras{молотарку} \textsuperscript{(47)}, яку крутить
\VUgras{локомобіль} \textsuperscript{(45)} через \VUgras{привідний пас}
\textsuperscript{(46)}, і так зерно відділяється від \VUgras{соломи}, не
покидаючи поля, де було посіяне.

%%K t08-tit-4 sub
\VUsaut{5.0mm}
\VUpk{10.2pt}{40}{Гроза.}
\VUsaut{3.4mm}

%%K t08-c1-12-1 p
12. --- Під час жнив часто зриваються сильні \VUgras{грози}
\textsuperscript{(53)}. Спершу здалеку чути гуркіт \VUgras{грому}; підіймається
сильний \VUgras{західний вітер} \textsuperscript{(54)} і до стогону гне
найбільші дерева \VUgras{лісу} \textsuperscript{(49)}. Чорні \VUgras{хмари}
\textsuperscript{(56)} беруть небо приступом; їх перетинають сліпучі зигзаги
\VUgras{блискавки} \textsuperscript{(55)}. Починає падати \VUgras{дощ}, а часом і
\VUgras{град}, прибиваючи готовий до жнив хліб.

%%K t08-c1-13-1 p
13. --- Часто блискавка запалює якусь будівлю. Так, торік дах \VUgras{вітряка}
\textsuperscript{(50)} згорів дощенту, а два його \VUgras{крила}
\textsuperscript{(51)} було розбито.

%%K t08-c1-14-1 p
14. --- За кілька годин вертається спокій. На обрії з'являється блискуча
\VUgras{веселка} \textsuperscript{(52)}, і природа знову береться за життя,
перерване на час. Селяни, що поховалися в \VUgras{куренях}
\textsuperscript{(48)}, вертаються до роботи і хоробро беруться направляти щойно
завдану шкоду.

%%K t08-tit-5 sub
\VUsaut{6.0mm}
\VUcentre{13.2pt}{90}{\textit{Друга сцена.}}
\VUsaut{3.6mm}

%%K t08-tit-6 sub
\VUpk{10.2pt}{40}{Полювання. --- Збір винограду.}
\VUsaut{1.4mm}
\VUpk{10.2pt}{40}{Риболовля.}
\VUsaut{4.0mm}

%%K t08-c2-01-1 p
1. --- Наприкінці \VUgras{серпня} або в середині \VUgras{вересня}, залежно від
краю, відкривається полювання. Ледве перші промені сонця виманять
\VUgras{ящірок} \textsuperscript{(2)} із нір, як \VUgras{мисливець}
\textsuperscript{(5)} уже виходить зі своїми \VUgras{собаками}
\textsuperscript{(4)}.

%%K t08-c2-02-1 p
2. --- Він надів свій міцний \VUgras{мисливський костюм} \textsuperscript{(9)} з
товстого сукна і шкіряні \VUgras{гетри}, які дозволять йому йти по мокрій
траві, де ховаються \VUgras{ропухи} \textsuperscript{(1)}; він узяв
\VUgras{рушницю} \textsuperscript{(6)} і \VUgras{ловецьку торбу}
\textsuperscript{(10)} --- і в дорогу.

%%K t08-c2-03-1 p
3. --- Запал полювання приводить його в саму середину виноградника, де збір у
повному розпалі. Діти виноградаря, які звичайно стережуть кіз коло дороги,
сьогодні пускають їх пастися де заманеться і, прив'язавши до кілка старого
\VUgras{цапа} \textsuperscript{(3)}, який напевно скористався б із волі, щоб
утекти, теж беруть свою пайку веселощів. Один дістає гроно з \VUgras{кошика на
виноград} \textsuperscript{(12)} і бавиться, вичавлюючи \VUgras{сік}
\textsuperscript{(14)} у \VUgras{чашку} \textsuperscript{(13)}, щоб вийшло сусло
(неперебродиле вино). Другий вінчає себе \VUgras{вінком із листя}
\textsuperscript{(15)}, який щойно сплів. Біля них зухвалі \VUgras{горобці}
\textsuperscript{(16)} підлітають до самого преса, щоб поцупити винограду.

%%K t08-c2-04-1 p
4. --- Поки діти бавляться, чоловіки працюють. \VUgras{Збирачі}
\textsuperscript{(18)} носять виноград до погреба в \VUgras{заплічних кошиках}
\textsuperscript{(17)}; \VUgras{бондар} \textsuperscript{(19)} лагодить
\VUgras{бочки} \textsuperscript{(20)}, перевіряє, чи цілі \VUgras{клепки}
\textsuperscript{(22)} і \VUgras{дна} \textsuperscript{(21)}, і заміняє луснуті
\VUgras{обручі} \textsuperscript{(23)}. Це він же зробив великі \VUgras{чани}
\textsuperscript{(25)}, куди зсипають \VUgras{виноград} \textsuperscript{(26)},
щоб він бродив.

%%K t08-c2-05-1 p
5. --- Коли бродіння скінчиться, вино зливають, а решту кладуть під
\VUgras{прес} \textsuperscript{(28)}; помалу закручуючи великий \VUgras{ґвинт}
\textsuperscript{(29)}, вичавлюють сік, який стікає в \VUgras{чан}
\textsuperscript{(32)} і дає ще \VUgras{вина} \textsuperscript{(31)}.

%%K t08-tit-8 sub
\VUsaut{6.0mm}
\VUpk{10.2pt}{40}{Пиво і сидр.}
\VUsaut{3.4mm}

%%K t08-c2-06-1 p
6. --- По стіні \VUgras{погреба} \textsuperscript{(27)} в'ється гілка
\VUgras{хмелю} \textsuperscript{(30)}. Тут це лише оздоба, але в інших краях, де
лоза дуже рідка, хміль розводять, щоб робити \VUgras{пиво}.

%%K t08-c2-07-1 p
7. --- Маленька \VUgras{яблуня} \textsuperscript{(33)} гнеться під яблуками, які
дадуть добірний \VUgras{сидр}, коли достигнуть. У своїй \VUgras{клітці}
\textsuperscript{(34)} \VUgras{канарка} \textsuperscript{(35)} і \VUgras{щиглик}
\textsuperscript{(36)} заздрісно дивляться на горобців, що пустують на волі.

%%K t08-c2-08-1 p
8. --- Скільки сягає око, по схилу пагорбів стоять ряди \VUgras{тичок}
\textsuperscript{(40)}, що підпирають чудові \VUgras{лози} \textsuperscript{(39)},
важкі від \VUgras{грон}.

%%K t08-c2-08-2 p
Цілий рік \VUgras{виноградар} \textsuperscript{(42)} працював, щоб ходити коло
свого \VUgras{виноградника} \textsuperscript{(38)}, і ось він винагороджений за
труди добрим урожаєм. \VUgras{Збиральниці} \textsuperscript{(41)} наповнюють
чани, поставлені на воза, якого тягнуть \VUgras{мули} \textsuperscript{(43)}, і
всі веселі.

%%K t08-tit-9 sub
\VUsaut{5.0mm}
\VUpk{10.2pt}{40}{Риболовля.}
\VUsaut{3.4mm}

%%K t08-c2-09-1 p
9. --- Те саме й у \VUgras{вудкарів} \textsuperscript{(44)}, які прийшли на
світанку розташуватися коло води зі своїми \VUgras{вудками}
\textsuperscript{(45)}.

%%K t08-c2-09-2 p
\VUgras{Риба} \textsuperscript{(47)} бере гачок, ледве вони опустять у воду свою
\VUgras{волосінь} \textsuperscript{(46)}, і вони ледве встигають змінити
\VUgras{принаду}, як нова жертва вже б'ється на кінці волосіні.

%%K t08-c2-09-3 p
Проте \VUgras{рибалка з сіттю} \textsuperscript{(48)}, який кидає свою сіть із
\VUgras{човна} \textsuperscript{(49)} на середину \VUgras{річки}
\textsuperscript{(50)}, ловить куди більше.

%%K t08-c2-10-1 p
10. --- Осінь --- пора добрих прогулянок: пішки, \VUgras{верхи}
\textsuperscript{(52)}, на \VUgras{велосипеді} \textsuperscript{(53)}, на
\VUgras{автомобілі} \textsuperscript{(54)}. \VUgras{Дороги} \textsuperscript{(51)}
чисті, і люди користаються з останніх погожих днів, бо \VUgras{ластівки}
\textsuperscript{(56)} уже збираються на дахах, щоб полетіти в теплі краї. Листя
старої \VUgras{ліщини} \textsuperscript{(57)} жовкне, \VUgras{горіхи} падають, а
високі \VUgras{дерева}, що стоять \VUgras{гаєм} \textsuperscript{(58)} на
\VUgras{підвищенні} \textsuperscript{(59)}, скоро оголяться.

%%K t08-c2-11-1 p
11. --- \VUgras{Сонце} \textsuperscript{(60)} встає пізніше, дні стають коротші,
\VUgras{зранку} починає чутися доволі різкий холодок, і вже зграї
\VUgras{слукв} \textsuperscript{(61)} покидають наш негостинний край.
\VUsaut{6.0mm}
\VUfilet{22mm}
